%%%%%%%%%%%%%%%%%%%%%%%%%%%%%%%%%%%%%%%%%
% The Legrand Orange Book
% LaTeX Template
% Version 2.3 (8/8/17)
%
% This template has been downloaded from:
% http://www.LaTeXTemplates.com
%
% Original author:
% Mathias Legrand (legrand.mathias@gmail.com) with modifications by:
% Vel (vel@latextemplates.com)
%
% License:
% CC BY-NC-SA 3.0 (http://creativecommons.org/licenses/by-nc-sa/3.0/)
%
% Compiling this template:
% This template uses biber for its bibliography and makeindex for its index.
% When you first open the template, compile it from the command line with the 
% commands below to make sure your LaTeX distribution is configured correctly:
%
% 1) pdflatex main
% 2) makeindex main.idx -s StyleInd.ist
% 3) biber main
% 4) pdflatex main x 2
%
% After this, when you wish to update the bibliography/index use the appropriate
% command above and make sure to compile with pdflatex several times 
% afterwards to propagate your changes to the document.
%
% This template also uses a number of packages which may need to be
% updated to the newest versions for the template to compile. It is strongly
% recommended you update your LaTeX distribution if you have any
% compilation errors.
%
% Important note:
% Chapter heading images should have a 2:1 width:height ratio,
% e.g. 920px width and 460px height.
%
%	Cities list:
%	New York
%	Beijing
%	London (city)
%	Singapore
%	Dubai
%	Tokyo
%	Los Angeles
%	Frankfurt
% 	Kuala Lumpur
%	Brussel
%	Milan
%	Berlin
%
%%%%%%%%%%%%%%%%%%%%%%%%%%%%%%%%%%%%%%%%%

%----------------------------------------------------------------------------------------
%	PACKAGES AND OTHER DOCUMENT CONFIGURATIONS
%----------------------------------------------------------------------------------------

\documentclass[11pt,fleqn,oneside]{book} % Default font size and left-justified equations

%----------------------------------------------------------------------------------------

\input{structure} % Insert the commands.tex file which contains the majority of the structure behind the template

\begin{document}

%----------------------------------------------------------------------------------------
%	TITLE PAGE
%----------------------------------------------------------------------------------------

\begingroup
\thispagestyle{empty}
\begin{tikzpicture}[remember picture,overlay]
    \node(background) at (current page.center) {\includegraphics[height=\paperheight]{background}};
\end{tikzpicture}
\endgroup

\newpage

% \begingroup
% \section{Abstract}\index{Abstract}
% \thispagestyle{empty}
% \addcontentsline{toc}{chapter}{\textcolor{trolleygrey}{Abstract}}
\usechapterimagefalse % If you don't want to include a chapter image, use this to toggle images off - it can be enabled later with \usechapterimagetrue
\chapter{Abstract}
\textbf{Ripa Exchange è un exchange ibrido-decentralizzato con il focus di abbassare il livello di ingresso per aprire 
nuove piattaforme di exchange ed al contempo offrire ai traders un ambiente sicuro ed efficiente per eseguire operazioni di 
trading giornaliere.}\\

Il Team di Ripa Exchange crede che, indipendentemente dagli sviluppi nel mondo delle criptomonete, è oggigiorno costoso per aprire,
gestire, e trovare fiducia nelle proprie operazioni di cambio valuta non solo per le risorse richieste per avviare
una piattaforma di scambio affidabile ma soprattutto per sviluppare la piattaforma in sè e trovare la liquidità 
necessaria per generare profitto nei primi 3-5 anni di attività.\\

Azione è richiesta ed azione è richiesta adesso. Gli utenti sono frustrati da exchange inaffidabili che chiudono scappando con i 
fondi, vengono hackerati oppure non sostengono il carico giornaliero di una'industria crescente come questa è.
Nonostante l'impegno degli imprenditori di exchange per offrire piattaforme di cambio valuta efficienti, affidabili e 
semplici da utilizzare i fondi necessari per costruire queste piattaforme sono nella fascia di 250-300 mila euro senza includere
il costo del personale per offrire supporto di livello platino agli utenti finali, l'infrastruttura di gestione, e costi giornalieri
dell'attività. Tali costi vengono sostenuti per avere una piattaforma discreta alla quale ci si lega ad unica via con una software
house per sviluppi futuri.\\

É obbiettivo di questo progetto di sviluppare una piattaforma di cambia valuta Open Source, efficiente, affidabile ed 
offrire liquidità\footnote{Grazie alla tecnologia RLSP (Ripa Liquidity Service Provider)} all'exchange così creato dal \textbf{primo}
giorno di attività in modo tale che gli imprenditori in questo settore possono focalizzarsi nel pubblicizzare la propria piattaforma,
offrire supporto tecnico di libello platino ai propri utenti e soddisfare tutte le eterogenee legislazioni nell'industria delle valute virtuali.
Mentre vogliamo che l'esperienza di trading per l'utente finale sia la migliore possibile offrendo un ambiente di trading sicuro ed affidabile.\\
\usechapterimagetrue
% \endgroup

%----------------------------------------------------------------------------------------
%	TABLE OF CONTENTS
%----------------------------------------------------------------------------------------

\chapterimage{chapter_head_1_Beijing.jpg} % Table of contents heading image
\renewcommand*\contentsname{Indice dei Contenuti}
\tableofcontents % Print the table of contents itself

%----------------------------------------------------------------------------------------
%	PART
%----------------------------------------------------------------------------------------


%\part{Part One}

%----------------------------------------------------------------------------------------
%	CHAPTER 2: Introduction
%----------------------------------------------------------------------------------------

\chapterimage{chapter_head_2_London.jpg} % Chapter heading image

\chapter{Introduzione}
Il livello di accesso all'industria delle valute virtuali è piuttosto elevato da un punto di vista tecnico per l'utente medio, 
ed ha un livello di accesso economicamente elevato per l'imprenditore che vuole avviare un'attività in questo campo, valore dato 
dall'acquisto del codice sorgente dell'exchange di criptovalute, dal reclutamento di personale operativo per le operazioni di DevOps,
per reclutare operatori di supporto tecnico, per sottostare alle leggi nazionali ed internazionali in materia di AML/KYC, per avere liquidità
dal primo giorno di avvio delle operazioni di cambio. Con questo progetto, Ripa Exchange, vuole abbassare questo livello di ingresso 
perchè \textbf{gestire un exchange è DIFFICILE} ed il Team Ripa Exchange vuole che il lettore e futuro imprenditore nel campo delle valute
virtuali si concentri/focalizzi su cose veramente importanti non cavilli burocratici che l'industria richiede perchè TU vuoi avviare un'attività
in questo campo e tu necessiti del codice sorgente per avviarla (ed il codice sorgente per avviarla ti deve essere concesso gratuitamente).

Per rafforzare/chiarire questo principio economico il punto è che per avviare un exchange è richiesto un investimento esiguo dal tuo 
venture capital di riferimento e anche con un tale investimento il ritorno dei fondi investiti all'angel investor non sono garantiti
nei primi 5 anni di attività.

Per avviare un servizio professionale di cambio valuta crediamo che il codice sorgente dell'exchange e la liquidità da offrire ai tuoi clienti
dal primo giorno di attività dovrebbero esserti forniti gratuitamente: non deve essere accettabile pagare \euro150.000,00 ad una software house
solo per avere una piattaforma che funzioni e per la quale dovete erogare altri \euro100.000,00-150.000,00 per brandizzarla, personalizzarla
per i vostri scopi legandovi in questo modo ad una compagnia/software house che potrebbe fallire in futuro e che vi farebbe ritrovare
senza poter fare ulteriori modifiche al vostro exchange in quanto non vi hanno mai consegnato il codice sorgente della piattaforma/del prodotto
su cui il vostro business fa affidamento/si basa.

\textbf{Noi crediamo che tutto questo debba essere gratuito}: dobbiamo offrirti la miglior tecnologia nel mercato attuale in modo tale che tu 
possa focalizzarti sulla tua attività mentre noi ci focalizziamo in sviluppare la tecnologia per gestire la tua attività/i tuoi affari
in maniera efficiente, sicura, responsiva e produttiva. Questo è il motivo per cui Ripa Exchange si prefigge lo scopo di sviluppare una rete di exchange
focalizzandosi/utilizzando un'architettura che è \textit{efficiente, sicura, responsiva, compliant e personalizzabile} in modo tale che ogni exchange
nella rete può basarsi su di fondamenta solide mentre è comunque possibile modificare la singola istanza secondo le personalizzazioni 
che la compagnia che gestisce quell'istanza richiede/quel singolo exchange richiede.

Per raggiungere questo scopo abbiamo scelto di sviluppare la tecnologia Ripa Liquidity Service Provider basata su ARK - a blockchain for consumer adoption - 
il cui obbiettivo primario è incrementare l'adozione delle tecnologie blockchain agli utenti finali focalizzandosi su due aree critiche: 
\underline{Una Efficiente e Sicura Tecnologia di Base} e \underline{Servizi Reali per Gente Comune}. L'ecosistema ARK è ancora ad uno stadio
iniziale di sviluppo: la futura implementazione di ARK 2.0 permetterà di eseguire smart contract nativi sulla blockchain, che permetteranno
alla blockchain ARK di competere con la sua concorrente Ethereum da un punto di vista tecnologico. RLSP permetterà a tutti gli exchange della rete Ripa
di condividere la stessa liquidità-lo stesso libro degli ordini in base ai parametri di condivisione della singola istanza dell'exchange Ripa
configurabili nella tua area amministratore.

Il Ripa Founder Team (RFT), come presentato in ripaex.io, opera nel nome del gruppo Ripa. L'RFT è responsaile per l'uso corretto dei finanziamenti
ottenuti nella campagna di scambio token RIPA TEC come spiegato di seguito in questo documento.

Il RFT comprende-impone che i risultati delle fasi di Prevendita e RIPA TEC saranno usati esclusivamente per il finanziamento del progetto 
\emph{Ripa Exchange} come spiegato in questo whitepaper, che saranno resi disponibili per estrazione nella piattaforma apposita resta disponibile
più tardi quest'anno tec.ripaex.io, e che dovranno risultare nella creazione di una entità legale che sarà chiamata \emph{Ripa Exchange}. La creazione di tale 
compagnia è pianificata per il primo quadrimestre del 2019.

Fino a tale data RFT opera invece di \emph{Ripa Exchange, una compagnia nel processo di essere incorporata}.


%------------------------------------------------

%This statement requires citation \cite{article_key}; this one is more specific \cite[162]{book_key}.

%------------------------------------------------

\section{Terminologia}\index{Terminologia}
\begin{description}
	\item[\textsc{Ripa Exchange:}] un exchange FIAT <-> CRYPTO sviluppato partendo dal codice sorgente di Peatio \cite{peatio}
	\item[\textsc{Ripa Blockchain:}] una blockchain DPOS in cui la liquidità ed il libro degli ordini sono condiviss tra tutti gli exchange della rete
	\item[\textsc{Ripa Token - XPX:}] un token crittograficamente sicuro tramite il protocollo DPOS scambiato sulla blockchain Ripa 
	\item[\textsc{RIPA:}] l'ecosistema finanziario DPOS composto da Ripa Exchange e Ripa Blockchain
	\item[\textsc{RIPAEX:}] il nome del progetto, sito web del progetto e domini del progetto
    \item[\textsc{RLSP:}] Ripa Liquidity Service Provider, un libro ordini condiviso tra gli exchange nella stessa rete Ripa
	\item[\textsc{ARK:}] una piattaforma per l'adozione consumer di tecnologie blockchain \cite{ark}
	\item[\textsc{ACES:}] Ark Contract Execution Services \cite{aces} fornisce semplici protocolli e strumenti per costruire robusti
	mercati di scambio tra blockchain basati sulla tecnologia ARK SmartBridge
	\item[\textsc{“,” or “.”:}] la notazione italiana per virgola decimale e punti di migliaia è stata scelta per la stesura di questo 
	whitepaper: detto questo “,” rappresenta un punto decimale e “.” rappresenta la separazione tra migliaia e multipli di migliaia 
\end{description}

%Lists are useful to present information in a concise and/or ordered way\footnote{Footnote example...}.
\pagebreak
\section{Roadmap}
Saranno percorse essenzialmente quattro fasi del progetto RipaEx:\\

\tcbset{roadmapBox/.style={colback=yellow!10!white,colframe=azure(colorwheel),
equal height group=nobefaf,width=(\linewidth-1pt), height from=4cm to 8cm, nobeforeafter,
	center title,
	valign=top, halign=left}}
\begin{center}
\begin{tcolorbox}[roadmapBox,
	title=\textbf{\textsc{Finanziamneto del progetto: XPX presale e RIPA TEC (WP2)}}]

	Questa fase riconosce l'esistenza di un interesse in questo sviluppo di mercato
	da un capo del Mondo all'altro in riguardo l'abbassamento del livello economico di ingresso
	per sviluppare exchange di cryptovalute.
	L'obbiettivo è di eseguire la prima analisi omnicomprensiva dello stato dell'arte per fondare
	le basi delle fasi sucessive del progetto e per sviluppare il primo prototipo
	funzionante di un exchange centralizzato con il codice di Peatio.\\
	\vspace{1cm}
	\centering\textbf{\textsc{Prevista fine della fase: Gennaio 2019}}.
\end{tcolorbox}

\resizebox{0.05\textwidth}{26pt}{$\Downarrow$}

\begin{tcolorbox}[roadmapBox,
	title=\textbf{\textsc{Apertura del primo exchange e sviluppo di strumenti e risorse (WP3)}}]

	La seconda fase prende i risultati della prima e sviluppa da qui
	una serie di strumenti e risorse che offrono guida concisa ed omnicomprensiva per 
	i player senza considerare lo Stato di incorporazione/sviluppo del progetto. Con la prima
	istanza di Ripa Exchange in esecuzione/aperta al pubblico i primi contatti con altri player economici
	nell'industria delle valute virtuali possono essere presi.\\
	\vspace{1cm}
	\centering\textbf{\textsc{Prevista fine della fase: Giugno 2019}}.
\end{tcolorbox}

\resizebox{0.05\textwidth}{26pt}{$\Downarrow$}

\begin{tcolorbox}[roadmapBox,
	title=\textbf{\textsc{Esecuzione/Disseminazione (WP 7/8) e Coordinamento del Progetto (WP1)}}]

	Durante tutta la durata del progetto,
	attività di disseminazione (WP 7/8) sono eseguite nelle quali i risultati dei singoli
	pacchetti di lavoro individuali sono disseminati ai relativi gruppi obbiettivo inclusi
	partners del progetto, supporters RipaEx, manager di exchange, partner bancari come anche
	tutti gli altri gruppi obbiettivo rilevanti al progetto RipaEx/in esecuzione.
	Questa fase copre una vasta scala di tecniche di disseminazione, da libri stampati ed 
	elettronici, a workshop e sessioni di training/hackatons, reti esistenti, tutti aventi 
	come obbiettivo di definire standard per la comunicazione tra exchange tra entità pubbliche e private.
	Un OVERARCHING pacchetto di lavoro riguardante la gestione del progetto dall'inizio alla fine, 
	assicurando coordinazione adeguata, quality assurance e controllo del budget (WP1).
\end{tcolorbox}

\resizebox{0.05\textwidth}{26pt}{$\Downarrow$}

\begin{tcolorbox}[roadmapBox,
	title=\textbf{\textsc{Sviluppo di un exchange ibrido-decentralizzato (WP 4-6)}}]

	Usando gli strumenti e le risorse sviluppate in WP3, il pacchetto di lavoro 4-6
	si focalizza mettere in pratica la conoscenza acquisita e gli strumenti. I tre
	pacchetti di lavoro riflettono tre punti focali (e gruppi target) nel network di 
	exchange creato per incorporare/stabilire demo di successo su scala locale, nazionale ed internazionale:
	l'incorporazione di istanze di Ripa Exchange locali (WP4), analisi tecnica della tecnologia
	Ripa Liquidity Service Provider (WP5), e primo MVP/POC dell'exchange ibrido decentralizzato (WP6).
	La fase demo dal cuore dell'azzione RipaEx; WP 2 e 3 focalizzano in produrre deliverables (consegne ad esempio
	strumenti) che abilitano attività di dimostrazione efficienti e di successo.
	\vspace{1cm}
	\centering\textbf{\textsc{Prevista fine della fase: Gennaio 2021}}.
\end{tcolorbox}
\end{center}

\section{RipaEx Partners - RipaEx Governance}\index{RipaEx Partners - RipaEx Governance}
La maggior parte dei parnter sono imprenditori nell'industria delle valute virtuali, tuttavia
un Istituto di Ricerca ed Organizzazioni Finanziari sono anche presenti. I partner sono:
\begin{description}
	\item[Coordinatore]: Ripa Exchanges Ltd
	\item[CoBeneficiari]: 
\end{description}
\subsection{RipaEx Governance}\index{RipaEx Governance}
La governance per la rete di exchange creati, lo sviluppo del codice sorgente, il possesso dei token XPX, la Fondazione Ripa.

\section{Sommario del Progetto}
\begin{enumerate}
	\item \textsc{cosa}: RipaEx è un progetto per facilitare la presa di standard per condividere liquidità tra mercati di asset
	criptosicuri (crypto asset marketplace).
	L'obbiettio di RipaEx è la promozione di codice sorgente condiviso per portafogli ed exchange nell'industria delle valute
	virtuali: è obbiettivo di questo documento di riferimento di dare informazioni dettagliate per futuri/prospective,
	sviluppatori di exchange oppure imprenditori in questo campo, per abilitare processi di decision-making corretti e per 
	assicurare il successo dei loro progetti proposti.
	IT SEEKS di analizzare il potenziale reale di applicazione nel Paese scelto di una rete di exchange e la sua posizione nel mercato.
	\item \textsc{cosa}: gli asset crittograficamente sicuri (aka: "crypto assets") sono un'alternativa agli asset emessi dalle borse
	valori dei vari Paesi che ne possiedono una. Anche se, tuttavia, molte borse valori mondiali danno la possibilità ai loro 
	utenti di verificare e gestire gli asset in loro possesso il processo non è sempre trasparente e semplice da fare e questo è il motivo 
	per cui a partire dal 2009 \cite{bitcoin} un nuovo tipo di asset (verificabile dalla comunità) è stato implementato
	per dare a piccoli, medi e grandi investitori completa trasparenza nella gestione dei loro asset di investimento.
	\item \textsc{quanto}: recenti sviluppo a livello dell'Unione Europea e globalmente stanno trasformando sia come le valute virtuali
	sono trattate sia il modo in cui le ICO (Initial Coin Offering) sono legislate. Questi sviluppi combinati hanno messo
	l'uso e la produzione di valute virtuali in una luce favorevole in costante ascesa. \\
	In Ottobre 2015 la Corte Europea di Giustiza ha sentenziato che bitcoin ed altre valute virtuali sono esenti dalla tassazione IVA. \\
	In Luglio 2016 la Commissione Europea ha adottato proposal per ammendare la legislazione vigente 4th Anti-Money Laundering Directive (4AMLD) 
	che introduca le figure degli exchange di valute virtuali ed offritori di wallet all'interno del framework europeo di adeguata verifica
	dell'utente ed antiriciclaggio \cite{EUAMLCrypto}.\\
	In Febbraio 2018 la Commissione Europea lancia l'Osservatorio EU sulle Blockchain (EU Blockchain Observatory and Forum) \cite{EUBOaF} 
	per osservare gli sviluppi chiave nella tecnologia blockchain, promuovere attori europei e rafforzare l'ingaggio europeo con
	diversi stakeholders coinvolti in attività correlate alla blockchain.
	\item \textsc{perchè}: c'è ancora poca regolamentazione sulle ICO e solo gli Stati Uniti d'America al momento
	hanno scritto una legislazione definendo i token delle ICO come securities \cite{SECICO}.
	\item \textsc{chi}: i risultati delle indagini di mercato indicano che il medio tech savvy dai 18 ai 45 anni di età è
	l'utente medio delle valute virtuali tuttavia compagnie in finanza stanno anche comincianto ad inserire schemi di 
	valute virtuali all'interno dei loro portafogli specialmente dalla presentazione del contratto bitcoin futures 
	da parte del Gruppo CME Group Inc. nella borsa valori di Chicago lo scorso 18 Dicembre 2017.
	\item \textsc{quanto}: la capitalizzazione di mercato totale dell'industria delle valute virtuali è stata stimata attorno
	ai 317 miliardi di USD \footnote{Dati Coinmarketcap Aprile 2018}
	ed è prevista la crecita a 5.000 miliardi di USD nel prossimo timeframe di 10 anni \cite{cryptoMCTenYears}.


	\item \textsc{where}: Local authorities are working with National Governments to make sure local exchangers
	in the national territory are complying with national and international AML/KYC regulations.
	Venture capitals and Angel Investors are starting to release financing solutions to start-ups 
	in the Fintech industry all over the world from America to Asia passing through Europe and some 
	Countries are starting state-owned cryptocurrencies schemes to test the exchange of goods \& services
	on those (distributed ledger) technologies \cite{petro}.
	\item \textsc{how much}: The average cost for starting your own crypto asses marketplace is around \$ 150,000.00 only for a running instance of your exchange platform:
	to that you need to add costs to customize the platform before launch and in the future, advertising your new business, running costs for servers,
	network operators, support center operators and legal department to comply with your State of incorporation AML/KYC legislations and general company laws.\\
	\textbf{That is the reason because we think owning the source code of your exchange software is the best way to run a business in this industry}.
	\item \textsc{how}: The main problems encountered in opening a FIAT <-> CRYPTO marketplace is to find trusted
	banking partners to comply with the many different AML/KYC rule and procedures to exchange virtual currencies
	to FIAT currencies.
	\item \textsc{what}: Classical types of exchanges operations are: 
		\begin{enumerate}[label*=\arabic*.]
			\item \textbf{one-way exchanges}: in which a centralized application has all the liquidity to offer to its potential users
			\item \textbf{two-way exchanges}: in which a centralized or decentralized platform match the selling requests with the buying requests
			of its users
		\end{enumerate}
		On this a sub-classification is also necessary:
		\begin{enumerate}[label*=\arabic*.]
			\item \textbf{FIAT <-> CRYPTO exchanges}: in which exchanges operations are performed between FIAT\footnote{Traditional central banks owned currencies like EUR, USD, GBP, JPY, others...} 
			currencies and virtual currencies
			\item \textbf{CRYPTO <-> CRYPTO exchanges}: in which exchanges operations are performed only between virtual currencies
		\end{enumerate}
	You can build a matrix based on the four configurations above to build the exchange operation platform of your needs.
	\item \textsc{with what}: The specifications to look when choosing for an exchange platform to run are:
		\begin{enumerate}[label*=\arabic*.]
			\item \textbf{code}: Open Source, Closed Source or hybrid solution
			\item \textbf{modularization}: separation between exchange engine (orders matching engine), UI and user registry
			\item \textbf{UI responsiveness}
			\item \textbf{compliance} with current industry standards
			\item \textbf{customization} of the exchange engine, trading currencies, UI and other aspects of the crypto asset marketplace platform...
			\item \textbf{security} of the funds: saving in cold wallets and hot wallets configurable
			\item \textbf{transparency} of the funds: proof of solvency of the exchange
			\item \textbf{Multi-Accounts trading}: easy to configure new virtual currency protocols
			\item \textbf{Multi-Accounts users}: possibility to interact with user accounts from Google, Facebook, Twitter 
			to login into the platform and FIDO Alliance security standards for personal credentials.
		\end{enumerate}	
	\textbf{Those are not only technical decisions to be made but also economical} especially the owning of the source code of your crypto asset
	marketplace platform is fundamental to make future customization of your exchange in an independent way compared to rely on a single
	software house that makes the customizations for you.
	\item \textsc{how}: Options for finding users for your exchanges operations are: targeted marketing campaigns, innovative features in the industry,
	fee level based on trading quantities, bonuses for first registration and trading quantities, affiliate marketing for paying users to take their friends
	to your exchange.
	\item \textsc{how}: For setting-up a crypto asset marketplace a project must take into account the following legislation:
		\begin{enumerate}[label*=\arabic*.]
			\item \textbf{AML/KYC}: \textit{Fourth Anti-Money Laundering Directive} if business set up in the European Union \cite{4AMLD} or the AML/KYC reference implementation
			to your crypto asset marketplace Country of incorporation (as an example \textit{Intelligence Reform \& Terrorism Prevention Act of 2004}
			written by FinCEN in the United States of America).\\
			International recommendations for undergoing AML/CFT verifications are given by the Financial Action Task Force on Money Laundering \cite{FATF}.
			\item \textbf{Payment Licence}: By far the biggest and most arduous task with regards to legitimising the 
			FIAT <-> CRYPTO exchanges operations is obtaining a \textit{PSD Licence} \cite{PSD}. 
			The PSD licence follows Council Directive 2007/64/EC and is applied in each country via its own national laws. 
			Costs of an IPPC licence can vary between \euro XXXX and \euro XXXX, depending on the size of operation.
		\end{enumerate}
	\item \textsc{who}: The nature of the business under consideration by the Ripa Exchange project (small scale,
	localised FIAT <-> CRYPTO exchanges operations), means that each enterprise likely to have 7 or 8 staff: N.2 developers, 
	N.1 network/security operator, N.1 administrative, N.2 client support operators, N.1 legal and tax advisor. \\
	The turnover of such an enterprise however, because of the high value of the end product, is likely to be more than 
	\euro 350,000 a year and could be several times higher. A business of
	this scale lends itself to the following possible company structures: A simple partnership;
	A limited company; A non-profit company or social enterprise; A worker co-operative.
	Financial Agencies are potential key actors, but the type of business they can set up will
	depend on their legal status which does vary from country to country.
	\item \textsc{how}: Potential sources of funds for a small-medium sized crypto asset marketplaces are: Bank Loans; Low
	Interest Loan Schemes; Commercial Credit; Equity financing; Business Angels venture
	capital. Having a robust Business Plan and financial guarantees are essential elements
	for securing funding. The European Investment Fund (EIF) of the EIB, offers support in
	the form of guarantees for SMEs.
	\item \textsc{why}: The arguments for crypto asset marketplaces are for financial freedom, decentralizing of the value-transferring operations, and
	owning for real your money. 
	There are other benefits, well documented, such as faster payments, long term gain based on deflationary economy and prediction of 
	Great Depressions like the one that hit the global economy in 2008. But
	above all, virtual currencies are the only direct competitor to centralized value-transferring operations done by central banks.

	\item \textsc{why}: There is consensus in the literature that the use of virtual currencies in place of fiat currencies will result in 
	higher financial freedom especially as they fit into the Austrian school of economy \cite{austrianTheory} 
	(TODO: add more on Austrian economics school)

	\item \textsc{why}: Benefits of virtual currencies schemes (TODO: put some numbers)

	\item \textsc{how}: Securing assets on the blockchains means basically performing three operations
		\begin{enumerate}[label*=\arabic*.]
			\item \textbf{Generating a random private key}
			\item \textbf{Converting the private key generated in (1) into a public key}: a common protocol making this conversion
			in the virtual currencies industry is the ECDSA curve algorithm
			\item \textbf{Converting the public key generated in (2) into a virtual currency address}: common protocols for making this conversion
			are hash functions SHA-256, Base58 encoding, Base32 encoding
		\end{enumerate}
	At this point any value sent to the virtual currency address generated in (3) is secured on the blockchain of choice and accessible
	only from the owner of the relative private key generated in (1).
	\item \textsc{where}: The two critical factors affecting the cryptocurrency industry are banks concurrence and State banning.
	Although a harmonisation throughout Europe would be beneficial to development of the industry both in terms of 
	taxation and warranty approvals, this is currently not the case. Each country has its specific legislation and tax
	regime for all exchanges operations involving FIAT money, and State banning is going to completely liberalization 
	of this activities like European Union to complete banning and imprisonment of operators in this industry like Bangladesh
	\cite{bitcoinLegality}. 

	\item \textsc{where}: The Asiatic countries of South Korea, China and Japan are the leader in the field of cryptocurrencies for number of 
	transactions for over 9 years with a proactive approach and favourable tax regime. At the beginning of 2017 in Japan bitcoin
	has been declared legal tender but China has recently declared illegal token sale and exchanges and local cryptocurrencies
	marketplaces are closing down.
	\item \textsc{where}: Any assessment of your local market should include: number of potential users to reach, 
	type of exchange to incorporate (FIAT <-> CRYTPO or CRYPTO <-> CRYPTO), type of virtual currencies protocol 
	to integrate (POW, DPOS, Masternodes, others...), types of services to offer (exchange only, advanced trading tools,
	payment processor, others...), if FIAT <-> CRYPTO exchange number of FIAT payments processors to accept (PayPal, OKPay,
	MoneyPolo, others...), number of others exchanges in your region.

	\item \textsc{who}: there are a number of options for dealing with Warranty/Customer protection issues: 
	creating consumer pressure by making clear to the end users that the possession of the private keys of their
	virtual currency addresses make \textbf{liable} for any loss of the private keys meaning nobody can help
	them recovering their funds if the their private keys are lost. Creating consumer pressure to not leave funds
	on exchanges ("\textit{Be Your own Bank!!}"), making them choose the licensed exchanges in the market. 
	\item \textsc{who}: While it is very expensive to insure money exchanges operations and money transmitting operations, 
	examples of customer protections in the industry are: Kraken platform which is offering Mt. Gox users partial refund
	of their lost, NEO community giving refund to the users involved in the BitGrail hacking, Ethereum supporters
	giving The DAO investors partial refunds, other hacking cases...

	\item \textsc{Recommendations for all/law compliance}: if you intend to incorporate a FIAT <-> CRYPTO exchange you should
	focus from the first instance on law compliance by studying the AML/KYC laws of the country of incorporation and
	finding bank partners to work with. Local financial Authority can help to comply with rules \& regulations and 
	local cryptocurrencies foundations can help you to tune your exchanges operations to perform targeted
	operations based on the customers interests in the country of incorporation.
	Promote cryptocurrency-friendly users in the area of interest.
\end{enumerate}


%----------------------------------------------------------------------------------------
%	CHAPTER 3: Ripa Exchange
%----------------------------------------------------------------------------------------

\chapterimage{chapter_head_3_Singapore.jpg} % Chapter heading image

\chapter{The Ripa Exchange}
Ripa Exchange is an crypto asset (fiat money or cryptocurrency or something) marketplace, following the latest 
industry standards and resting in the principles of "\textit{open source, secure and efficient}". Ripa Exchange
aims to serve a platform for crypto-currency enthusiasts by providing a safe, secure, UI responsive, customizable 
and easy to use exchange that embraces open source principles and public trust.\\\\
Ripa Exchange is implemented with the Rails framework and other cutting-edge technology and will be migrated to an
hybrid-decentralized exchange where all exchanges in the Ripa network will share liquidity thanks to the RLSP technology.

\section{Mission}
\begin{quotation}
	``\textit{Our mission is to build the world best open source crypto asset marketplace with a high performance trading engine 
	and safety which can be trusted and enjoyed by users. Additionally we want to move the crypto currency exchange technology 
	forward by providing support and add new features. We are helping people to build easy their own exchange around the world.}''
\end{quotation}

Help is greatly appreciated, feel free to submit pull-requests or open issues.

\section{Features}
A free, transparent and internationalized open source crypto currency exchange.\\

\tcbset{featureBox/.style={colback=yellow!10!white,colframe=white!20!black,
equal height group=nobefaf,width=(\linewidth-1pt)/3, height=4.5cm, nobeforeafter,
	center title,
	valign=top, halign=left}}
\begin{tcolorbox}[featureBox,
	title=\textsc{Open Source}]

	\small	All source code are fully released under the terms of the MIT License.\\\vspace{5mm}
	\tiny Ripa Exchange is a customizable cryptocurrency exchange solution architecture enables easy connection to KYC/AML, 
	authentication, ETL/reporting, and other services.
\end{tcolorbox}
\begin{tcolorbox}[featureBox,
	title=\textsc{Compliant}]

	\small	International KYC/AML standards.\\\vspace{5mm}
	\tiny Ripa Exchange KYC efficiently submits and exchanges KYC information 
	to meet the banking supervisory standards and comply with Customer Due Diligence (CDD) requirements.
\end{tcolorbox}
\begin{tcolorbox}[featureBox,
	title=\textsc{Transparent \& Configurable}]

	\small	Customize in your own way\\\vspace{5mm}
	\tiny Major functions have been embedded in the source code – neat registration and log-in interface, 
	personalized deposit and withdraw procedure, best match of bid and ask, etc. These functions are comprehensive 
	and are ready to use with no extra work needed. 
\end{tcolorbox}

\begin{tcolorbox}[featureBox,
	title=\textsc{Internationalization}]

	\small	All users are able to view Ripa Exchange in a language to their best convenience.\\\vspace{5mm}
	\tiny Supporting many common languages, Ripa Exchange makes it easy for users to operate in their mother tongue. 
	You are encouraged to contribute to our language variety. Users will benefit from your efforts.
\end{tcolorbox}
\begin{tcolorbox}[featureBox,
	title=\textsc{Proof of Solvency}]

	\small	Easy deployable PoS.\\\vspace{5mm}
	\tiny Ripa Exchange Proof of Solvency (PoS) allows users to verify 
	the solvency of the Ripa Exchange based cryptocurrency exchange without compromising user privacy.
\end{tcolorbox}
\begin{tcolorbox}[featureBox,
	title=\textsc{Multi-Accounts Trading}]

	\small	Easy currency configuration.\\\vspace{5mm}
	\tiny Ripa Exchange allows to create multiple accounts and trading in multiple currencies. 
	Ripa Exchange makes it is easy to trade different currencies.
\end{tcolorbox}

\begin{tcolorbox}[featureBox,
	title=\textsc{Multi-Accounts Users}]

	\small	Easy account configuration.\\\vspace{5mm}
	\tiny Ripa Exchange allows to create multiple login accounts Google, Facebook, Twitter and FIDO Alliance
	login standards to secure your account.
\end{tcolorbox}
\begin{tcolorbox}[featureBox,
	title=\textsc{Enterprise Exchange}]

	\small	Start small, grow big.\\\vspace{5mm}
	\tiny Ripa Exchange enterprise exchange features include a high-performance matching engine, 
	scalable distributed worker threads, and SMS 2-factor authentication.
\end{tcolorbox}
\begin{tcolorbox}[featureBox,
	title=\textsc{Functional \& Intuitive}]

	\small	For the new trader, for the experienced trader.\\\vspace{5mm}
	\tiny Clean, user friendly registration and login interface. 
	Personalized deposit and withdraw procedure and a built-in proof-of-solvency audit.
\end{tcolorbox}

\section{Functional Analysis}
\begin{enumerate}
	\item \textsc{Why Redis-RabbitMQ is needed}:
	\item \textsc{Why NodeJS}: 
	\item \textsc{Why Pusher}:
	\item \textsc{ER MySQL}:
	\item \textsc{Ruby folders hierarchy}:
\end{enumerate}

\section{Technology Stack}
\begin{description}
	\item[\textsc{Ruby on Rails}]:
	\item[\textsc{MySQL}]: 
	\item[\textsc{Redis}]:
	\item[\textsc{RabbitMQ}]:
	\item[\textsc{NodeJS}]:
	\item[\textsc{Pusher}]:
\end{description}

\section{User Interface}
Ripa Exchange user interface is based on Peatio user interface a UI responsive user interface built in Ruby on Rails and 
completely separated from the exchange order engine.
Design for a customised UI are in place to offer to Ripa Exchange users the best experience on all devices: here you can find
some screenshot of the current exchange interface design.

\subsection{End-User Interface}
Following some screenshot of the end-user Ripa Exchange-Peatio interface: 
\textcolor{darkred}{those screenshots are a work in progress and they may not represents the user interface of the final product}\\
\begin{center}
	\includegraphics[width=\textwidth*2/3,height=\textheight,keepaspectratio]{RE_TradingUIC}
	\captionof{figure}{Ripa Exchange trading UI}
	\includegraphics[width=\textwidth*2/3,height=\textheight,keepaspectratio]{RE_depositBTCC}
	\includegraphics[width=\textwidth*2/3,height=\textheight,keepaspectratio]{RE_historyC}
	\includegraphics[width=\textwidth*2/3,height=\textheight,keepaspectratio]{RE_solvencyC}
	\captionof{figure}{Ripa Exchange deposit/whitdraw, order history and solvency screen}
\end{center}

\subsection{Admin Interface}
Following some screenshot of the administrative console of Ripa Exchange-Peatio interface: 
\textcolor{darkred}{those screenshots are a work in progress and they may not rapresents the user interface of the final product}\\
\begin{center}
	\includegraphics[width=\textwidth*2/3,height=\textheight,keepaspectratio]{RE_adminDashboardC}
	\includegraphics[width=\textwidth*2/3,height=\textheight,keepaspectratio]{RE_adminDepositsC}
	\includegraphics[width=\textwidth*2/3,height=\textheight,keepaspectratio]{RE_adminWithdrawsC}
	\captionof{figure}{Ripa Exchange Admin console dashboard, deposit/whitdraw screen}
	\includegraphics[width=\textwidth*2/3,height=\textheight,keepaspectratio]{RE_adminUserInfo1C}
	\includegraphics[width=\textwidth*2/3,height=\textheight,keepaspectratio]{RE_adminVerifyAccountC}
	\captionof{figure}{Ripa Exchange Admin user info and verify account screen}
\end{center}

\section{Features at Launch}
\begin{description}
	\item[\textsc{CRYPTO <-> CRYPTO}]:
	\item[\textsc{OAuth}]: Facebook, Google, Twitter
	\item[\textsc{FIDO}]: login with FIDO Alliance standards
	\item[\textsc{Protocols}]: POW, DPOS, Masternode, tether, ERC20
	\item[\textsc{Currencies}]: BTC, ETH, DOGE, BCH, TUSD, ARK, LISK, SHIFT, RISE, KAPU, OXY, RIPA, promising ERC20 tokens
	\item[\textsc{Main Markets}]: BTC, ETH, ARK
	\item[\textsc{Orders Type}]: market, limit
\end{description}

\subsection{Future Features}
\begin{description}
	\item[\textsc{E-Wallets}]: OKPay, NETELLER, MoneyPolo, others...
	\item[\textsc{Advanced Trading Features}]: margin trading, stop loss and take profit, 
	\item[\textsc{FIAT <-> CRYPTO}]
	\item[\textsc{Other tools}]: VISA/MasterCard, merchants tools, P2P Lending 
\end{description}


\section{Towards a Decentralized Exchange...}
Ripa Exchange is a centralized exchange which will be converted into an hybrid-decentralized exchange to create a network of exchanges
that share the same liquidity between each of them so you can offer liquidity from day 1 of your exchanges operations.\\

To offer the PROs of a centralized exchange with all the PROs of a decentralized exchange without the CONs of both of them, 
we need to convert first into an hybrid-decentralized exchange and during phase 3 of our project (WP4-6) we will make all the 
functional and technical analysis required to make the next step on solid ground.

%----------------------------------------------------------------------------------------
%	PART
%----------------------------------------------------------------------------------------

%\part{Part Two}

%----------------------------------------------------------------------------------------
%	CHAPTER 4: Ripa Blockchain
%----------------------------------------------------------------------------------------

\chapterimage{chapter_head_4_Dubai.jpg} % Chapter heading image
\chapter{The Ripa Blockchain}
The project RipaEx will have its own blockchain called Ripa which will be run on the DPOS protocol and has the \PHP 
token (XPX symbol) running on it that will serve the five following purposes:
	\begin{enumerate}
		\item to list new cryptocurrencies on Ripa Exchanges
		\item to advertise new projects
		\item to buy RipaEx gadget on Ripa Exchange Store
		\item to pay for the sell of goods \& services on authorized resellers with our RipaEx POS (Point of Sale)
		\item to share liquidity between Ripa Exchanges in the same network
	\end{enumerate}

Ripa \PHP or XPX is a cryptocurrency derivate from ARK, Lisk, Crypti and BitShares with unique differences and improvements
for reaching the goal of shared liquidity between exchanges in the same Ripa network. This code however inherits the 
simplified interactions between ARK and other blockchain systems using DPoS as their consensus. This homogeneous codebase
allows for the potential to provide service bridges in the form of ARK-Lisk blockchain apps, along with any other additional 
systems provided by their blockchain administrators.

\vspace{5mm}
\textsc{\textbf{Ripa blockchain is an ARK fork and the use of a blockchain technology for the network of exchanges created will complete
the ripa ecosystem by permitting each exchange in the ripa network to share the same liquidity. We will always entrust ARK as our 
blockchain technology provider to merge their code into our latest features for what concerning the RIPA Blockchain}}

\vspace{5mm}

Explained the use of the Ripa Blockchain in the RipaEx ecosystem and explained the RIPA - ARK technological relations what is 
following are the specifications of the Ripa Blockchain derived/inherited from ARK.

\section{Delegated Proof of Stake Technology}
Ripa Blockchain, will inherits the Delegated Proof of Stake (DPoS) consensus system that was first introduced
by BitShares. This consensus algorithm was designed to eliminate the issues associated with Proof of Work (PoW), 
namely the centralization of computing power and the exponentially increasing waste of real world energy. While not
completely decentralized as it relies on consensus by a fixed number of elected delegates, it guarantees a better decentralization 
than Bitcoin. The consensus algorithm implementation is improved over time, evolving into an optimal consensus system.\\

The technical description of the Ripa blockchain is as follow:
\begin{enumerate}
	\item DPoS (Delegated Proof of Stake)
	\begin{description}
		\item - 101 active forging Delegates
		\item - Delegates selected by vote mechanism built into DPoS
		\item - 115,000,000 RIPA - Seeded Genesis Block
	\end{description}
	\item Multi-signature accounts
	\item Constant block reward
	\begin{description}
		\item - 2 \PHP per block
		\item - Inflation Rate (with 8s block times)
		\begin{description}
			\item 6.31\% for the first year
			\item 5.93\% the 2nd year
			\item 4.02\% the 10th year
		\end{description}
		\item - 8-second block time
		\begin{description}
			\item - Decreased block time possible with future upgrades to the core.
		\end{description}
		\item - 25 transactions per block
		\begin{description}
			\item - Increased via soft fork as needed.
		\end{description}
	\end{description}
	\item Routing tables
	\item SmartBridge data field for custom use and bridging blockchains (ARK Contract Execution Service)
	\item Batch transactions\footnote{Future upgrade: when ARK 2.0 will be released \label{note1}}
	\item Custom transaction fees\footnotemark[\value{footnote}]
	\item Native smart contract execution\footnote{Future upgrade: when ARK Virtual Machine will be released}
\end{enumerate}

\vspace{5mm}
\textsc{\textbf{As said Ripa blockchain is an ARK fork and we will entrust ARK as our blockchain technology provider to merge their 
code into our latest features for what concerning the RIPA Blockchain}}, features like:
\begin{description}
	\item[Scaling network] up to the level of major Credit Card networks with core upgrades
	\begin{description}
		\item - Increasing the number of Forging Delegates
		\item - Increasing the Block Size to include more transactions
		\item - Implementation of pre-approval PBFT block concept testnet [codename:
		TwinChain]
		\item - Routing tables, to minimize hops among nodes when blocks are broadcast
		\item - Include forging with RIPA Uncles.	
	\end{description}
	\item[Two node types] is used to secure the RIPA network:
	\begin{description}
		\item[Relay nodes] - Nodes with full API functionality, acting as a backend for the
		feature rich lite clients. Relay nodes do not collect any transaction fee and do
		not have the ability to Forge RIPA Blocks.
		\item[Forging nodes] - Nodes with reduced API functionality, decreasing the
		exposure to potential DDoS attacks on the RIPA Platform. Forging nodes are
		able to Forge RIPA and receive transaction fees.
	\end{description}
	\item[Official lite client] for network access is be provided shortly before the mainnet
	launch including desktop clients (Windows, MacOS, and Linux) and mobile clients
	(Android and iOS).
	\item[OffLine wallet creation]: the network itself does not use a Graphical User Interface by default. Any RIPA
	account can be created offline and managed at no cost with a single device
	(computer, mobile phone, embedded ARM, IoT).
\end{description}

\section{Hierarchical Deterministic (HD) Wallets (BIP32)}
The structure of the public and private key generation follows the same specification
as Bitcoin. A custom implementation of BIP32 for Hierarchical Deterministic Wallets
is provided to RIPA users.

\section{Fees}
The fee for standard transactions is set at \PHP0.1 but the will be flexible in future releases of Ripa Blockchain.
At mainnet launch, a fee structure is provided out of the box to forging delegates with the following rules:
	\begin{description}
		\item - Transaction \PHP0.1
		\item - Vote \PHP1 (101 votes per transaction)
		\item - Second Signature \PHP1
		\item - Multi Signature \PHP1 per signature + \PHP1 per signing account
		\item - Registering a delegate \PHP25.
	\end{description}
All fees are paid to the forging node which processes the block containing those fees.

\section{Ripa Delegates and Delegate Voting}
Any node running the core blockchain code wishing to become a forging node must
register their account within the RIPA network. The fee for this registration is set to
25 \PHP per delegate account registered.
RIPA incorporates a new DPoS voting system originally envisioned by the Crypti
Founders. The RIPA system fee is 1 \PHP per delegate vote. The voting weight of each
wallet will be split evenly between all delegates voted. For example:
\begin{description}
	\item - If a wallet votes for one delegate, that delegate receives 100\% of the wallets
voting weight.
	\item - If that wallet votes for an additional delegate, the entire vote weight is split
evenly between the two delegates at 50\%.
	\item - By adding a third delegate, the voting weight splits again, and each of the
three delegates receives 33.333\% of the voting weight from that wallet.
\end{description}

The 101 forging nodes with the highest number of votes are eligible to Forge RIPA
blocks. This design eliminates the possibility that any single large RIPA holder or an
organization holding large percentages of RIPA are able to gain control over the
entire network by voting for all of their nodes into forging positions, thus effectively
taking complete control over that DPoS Blockchain. Votes from RIPA Tokens held by
RIPA Crew may be used at RIPA Crew's discretion.

\section{Connecting Blockchains}
\subsection{ARK Bridged Blockchains (SmartBridge Technology)}
The ARK Platform does not provide direct support for sidechains or dapp databases.
Instead, a mechanism to bridge together blockchains is provided via a bridging
function built into ARK Core where any blockchain can send and receive trigger
function notices and informational data through the primary ARK network via
custom developed \textbf{SmartBridge(s)} and \textbf{Encoded Listeners}.
\begin{center}
	\includegraphics[width=\textwidth,height=\textheight,keepaspectratio]{ARKSmartBridge}
	\captionof{figure}{ARK SmartBridge Technology}
\end{center}

\subsection{A.C.E.S. - ARK Contract Execution Service}
ACES is a blockchain interoperability platform that provides simple protocols and tools for building a robust blockchain 
service marketplace.

ACES is composed mainly of the following three components:
\begin{description}
	\item[Listeners] ACES Listeners provide a way for all the different blockchain transaction events to be easily 
	consumed via a common REST-ful API. The API allows consumers to create subscriptions and receive blockchain events 
	in real-time using Webhook callbacks.
	\item[Services] ACES Services create and excute Service Contracts, which can be anything from uploading 
	a file to a storage blockchain, performing value transfers, creating smart contracts, executing code on 
	blockchain based computing platforms, or interacting with IoT hardware.
	\item[Marketplace Console] The ACES Marketplace Console is a consumer dashboard for searching and executing 
	service contracts listed on the Marketplace. ACES Service providers can list their service nodes using the Marketplace API.
\end{description}

\begin{center}
	\includegraphics[width=\textwidth,height=\textheight,keepaspectratio]{ACES}
	\captionof{figure}{ARK-BTC A.C.E.S. SmartBridge Implementation Overview}
\end{center}

\section{Ripa Liquidity Service Provider (R.L.S.P)}
Using the SmartBridge technology RipaEx will build a mechanism to share liquidity between all the exchanges in the Ripa network
by writing the single exchange orderbook in the Ripa Blockchain and by executing order matching between all exchanges in the network.

In this way you can have the benefits of a decentralized orderbook (like liquidity) with the benefits of a centralized exchange (like
platinum customer support and FIAT exchange).

\section{Ripa Community Fund}
To permits the born of new exchanges in the Ripa network a Ripa Community Fund - RCF - is created with the following characteristics:
\begin{description}
	\item[Starting Principal]: the RCF will have a starting operating capital of 5\% (5,750,000 \PHP) of the genesis block
	\item[Recurring Participation]: each delegate will contribute to the RCF with 5\% of its forged XPXs
\end{description}

To obtain funds from the RCF to start your own Ripa Exchange you will submit your proposal in the official RCF section in the Ripa forum at
anytime after the first Ripa Exchange instance will be operative.

%----------------------------------------------------------------------------------------
%	CHAPTER 5: Token Sale
%----------------------------------------------------------------------------------------

\chapterimage{chapter_head_5_Tokyo.jpg} % Chapter heading image
\chapter{Token Sale}
\section{Introduction}
The RipaEx XPX token sale will be separated into two phases:
	\begin{description}
		\item[\textsc{PreSale}]: that will run from April to September 2018
		\item[\textsc{RIPA TEC}]: that will run from October to December 2018
	\end{description}

\subsection{Exchange Rate \& Accepted Cryptocurrencies}
During the token sale the exchange rate for the XPX tokens will be \textbf{\PHP/\euro0.10}.\\

The following cryptocurrencies are accepted: \textbf{BTC, ETH, ARK, LISK}.\\

If you want to send us a different cryptocurrency send us your choice and we will tell you immediately if we can make 
the exchange and how to execute the operation in detail (receiving address, exchange rate, quantity)\\

\textsc{Exchange rate is Bitstamp last and is valid 60 minutes from quote}.

\section{How to Invest}
\subsection{XPX Ripa Token PreSale}
Private sale of XPX tokens has \textbf{started on Monday 04/02/2018 at 00:00 GMT and will end on Sunday 09/30/2018 at 23:59 GMT}.

You can buy directly here: 
\begin{center}
	\href{https://www.ripaex.io/privsale}{\textsc{www.ripaex.io/privsale}}
\end{center}

or send us your buying enquiries on our social Slack, Telegram, Facebook, Twitter or on other social networks where we are present.
For all the duration of the private sale period every exchange operation \textbf{will receive a 100\% bonus on its sending amount
until the first 5.000.000 XPX tokens} have been exchanged and will receive a 50\% bonus thereafter.

\subsection{RIPA TEC}
Fundraising through the RIPA TEC platform will follow the schedule:
\begin{description}
	\item[from 10/01/2018 00:00 to 10/14/2018 23:59]: with bonus of 50\% 
	\item[from 10/15/2018 00:00 to 10/28/2018 23:59]: with bonus of 40\% 
	\item[from 10/29/2018 00:00 to 11/11/2018 23:59]: with bonus of 30\% 
	\item[from 11/12/2018 00:00 to 11/25/2018 23:59]: with bonus of 20\% 
	\item[from 11/26/2018 00:00 to 12/09/2018 23:59]: with bonus of 10\% 
	\item[from 12/10/2018 00:00 to 12/23/2018 23:59]: with bonus of 5\% 
\end{description}
Coins distribution will be available at the end of each trading day and performed automatically through the RIPA TEC platform. Minimum
threshold to reach is 25 BTC.\\

\textsc{As the XPX exchange rate is calculated against the euro currency directly you are guaranteed to not loose any invested fund
during the sale period as the XPX tokens will be listed on our exchanges at the beginning of 2019 directly with the same \textbf{\PHP/\euro0.10}
exchange rate agreed during the token sale-ICO phase}.

\section{XPX Ripa Token Distribution}
\textbf{115,000,000} XPX tokens are seeded into the genesis block: distribution of the XPX tokens generated are described into
figure \ref{fig:distribution} and table \ref{tab:distribution}.

\vspace{5mm}
%\begin{center}
	\begin{tikzpicture}
		\pie [rotate = 180, explode={0.2, 0.1, 0.1, 0.1, 0.1, 0.1}, radius = 3]
		{65/PreSale \& RIPA TEC,
		15/RIPA Founders Team, 
		6/for marketing,
		5/for bounties,
		5/for RCF,
		4/ARK Team}
	\end{tikzpicture}
	\captionof{figure}{Ripa token distribution - Division of funds}	
	\label{fig:distribution}
%\end{center}

\vspace{5mm}
\begin{table}[H]
	\centering
	\begin{tabular}{l l l}
		\toprule
		\textbf{Percentage (\%)} & \textbf{Quantity (\PHP)} & \textbf{Purpose} \\
		\midrule
		65		& 74,750,000	& to distribute in PreSale and RIPA TEC	\\
		15      & 17,250,000	& to the RIPA Founders Team	\\
		6       & 6,900,000		& for marketing	\\
		5       & 5,750,000 	& for bounties	\\
		5       & 5,750,000		& for Ripa Community Fund - RCF	\\
		4       & 4,600,000		& to the ARK Team	\\
		\bottomrule
	\end{tabular}
	\captionof{table}{Ripa token distribution}
	\label{tab:distribution}
\end{table}

\vspace{5mm}
\textbf{Tokens not sold during the phases of PreSale and RIPA TEC will be burnt} forever to permits and equal distributions 
of the funds and avoid speculations on the token remaining.

\section{Division of Funds}
The funds collected during the PreSale and RIPA TEC phases will be divided as explained in figure \ref{fig:division} and 
table \ref{tab:division} below

\vspace{5mm}
\begin{tikzpicture}
	\pie [rotate = 180, explode={0.2, 0.1, 0.1}, radius = 3]
	{60/for project development,
	20/for marketing, 
	20/for legal support}
\end{tikzpicture}
\captionof{figure}{Ripa division of funds}	
\label{fig:division}

\vspace{5mm}
\begin{table}[H]
	\centering
	\begin{tabular}{l l}
		\toprule
		\textbf{Percentage (\%)} & \textbf{Purpose} \\
		\midrule
		60		& for project development	\\
		20		& for marketing	\\
		20		& for legal support	\\
		\bottomrule
	\end{tabular}
	\captionof{table}{Division of funds}
	\label{tab:division}
\end{table}

\vspace{5mm}
where the row project development funds allocation is described in table \ref{tab:focus} below.

\vspace{5mm}
\begin{table}[H]
	\centering
	\begin{tabular}{l l}
		\toprule
		\textbf{Percentage (\%)} & \textbf{Purpose} \\
		\midrule
		55		& functional analysis, technical analysis, development, testing	\\
		15		& infrastructure	\\
		10		& security	\\
		20		& technical support	\\
		\bottomrule
	\end{tabular}
	\captionof{table}{Project development focus}
	\label{tab:focus}
\end{table}

%----------------------------------------------------------------------------------------
%	CHAPTER 6: Economic Projections
%----------------------------------------------------------------------------------------

\chapterimage{chapter_head_6_Frankfurt.jpg} % Chapter heading image

\chapter{Business Model}

\section{Project Milestones}
A free, transparent and internationalized open source crypto currency exchange.\\

\begin{tcbraster}[raster columns=5,raster rows=1,raster height=0.8cm,
	valign=center, halign=center,
	enhanced,size=small,sharp corners,colframe=azure(colorwheel),coltext=white,
	colback=azure(colorwheel),fit algorithm=hybrid* ]
	\tcboxfit{}
	\tcboxfit{\textbf{25 BTC}}
	\tcboxfit{\textbf{250 BTC}}
	\tcboxfit{\textbf{1000 BTC}}
	\tcboxfit{\textbf{2000 BTC}}
\end{tcbraster}
\begin{tcbraster}[raster columns=5,raster rows=4,raster height=10cm,
	valign=center, halign=center,
	enhanced,size=small,sharp corners,colframe=silver,coltext=black,
	colback=silver,fit algorithm=hybrid* ]
	\tcboxfit{\textsc{\textbf{Development}}}
	\tcboxfit{\tcbfontsize{0.8} CRYPTO <-> CRYPTO exchange release with 25 cryptocurrencies supported and 3 main trading pairs}
	\tcboxfit{\tcbfontsize{0.8} CRYPTO <-> FIAT exchange release with MasterCard prepaid card. 50 cryptocurrencies supported and 3 main trading pairs}
	\tcboxfit{\tcbfontsize{0.8} Advanced trading features. Open of a second and third Ripa Exchanges to join the Ripa network}
	\tcboxfit{\tcbfontsize{0.8} Open of ten Ripa Exchanges across the Globe}

	\tcboxfit{\textsc{\textbf{Marketing}}}
	\tcboxfit{\tcbfontsize{0.8} Support from one agency, adv campaigns on targeted channels}
	\tcboxfit{\tcbfontsize{0.8} Support from three agencies, adv capaigns pushing harder, store with RipaEx gadgets}
	\tcboxfit{\tcbfontsize{0.8} International presentation event in London}
	\tcboxfit{\tcbfontsize{0.8} Support from five agencies, targeting all UN countries, presentation event in Tokyo}

	\tcboxfit{\textsc{\textbf{Legal}}}
	\tcboxfit{\tcbfontsize{0.8} Legal department working on AML/KYC international standards}
	\tcboxfit{\tcbfontsize{0.8} Legal departments working on AML/KYC international standards in New York, London, Hong Kong}
	\tcboxfit{\tcbfontsize{0.8} International Government Agencies cooperation}
	\tcboxfit{\tcbfontsize{0.8} Banking parnership for defining AML/KYC standards}

	\tcboxfit{\textsc{\textbf{Blockchain}}}
	\tcboxfit{\tcbfontsize{0.8} DevNET development for RLSP and future ARK releases}
	\tcboxfit{\tcbfontsize{0.8} Contributions to ARK development for AVM}
	\tcboxfit{\tcbfontsize{0.8} Contributions to ARK development for standard timing releases}
	\tcboxfit{\tcbfontsize{0.8} Partnership with ARK for blockchain technology developments}
\end{tcbraster}

\section{Market Overview}

\section{Local Market Analysis}

\section{5 Years Projections}

\section{XPX Token Economics}

%----------------------------------------------------------------------------------------
%	CHAPTER 7: Conclusion
%----------------------------------------------------------------------------------------

\chapterimage{chapter_head_7_LosAngeles.jpg} % Chapter heading image

\chapter{Team \& Conclusion}

\section{Team}
\begin{tcbraster}[raster columns=1,raster rows=1,raster height=0.4cm,
	valign=center, halign=center,
	enhanced,size=small,sharp corners,colframe=azure(colorwheel),coltext=white,
	colback=azure(colorwheel),fit algorithm=hybrid* ]
	\tcboxfit{}
\end{tcbraster}
\begin{tcbraster}[raster columns=2,raster rows=4,raster height=10cm,
	valign=center, halign=left,
	enhanced,size=small,sharp corners,colframe=silver,coltext=black,
	colback=silver,fit algorithm=hybrid* ]
	\tcboxfit{\tcbfontsize{0.8} \textbf{Giovanni Silvestri @ BitNow}\\
	Founder and Head of the Project\\
	10+ Years experience in IT Field,\\
	4+ Years in Cryptocurrencies Consultancy and Private Exchanging,\\
	2+ Years Experience in finance\\
	Location: Treviso (TV), Italy\\
	Bitcointalk
	Linkedin}
	\tcboxfit{\tcbfontsize{0.8} CRYPTO <-> FIAT exchange release with MasterCard prepaid card. 50 cryptocurrencies supported and 3 main trading pairs}

	\tcboxfit{\tcbfontsize{0.8} Support from three agencies, adv capaigns pushing harder, store with RipaEx gadgets}
	\tcboxfit{\tcbfontsize{0.8} International presentation event in London}

	\tcboxfit{\tcbfontsize{0.8} International Government Agencies cooperation}
	\tcboxfit{\tcbfontsize{0.8} Banking parnership for defining AML/KYC standards}

	\tcboxfit{\tcbfontsize{0.8} Contributions to ARK development for standard timing releases}
	\tcboxfit{\tcbfontsize{0.8} Partnership with ARK for blockchain technology developments}
\end{tcbraster}


\section{Recommendations}

\section{Social}

%----------------------------------------------------------------------------------------
%	CHAPTER 8: Legal
%----------------------------------------------------------------------------------------

\chapterimage{chapter_head_8_KualaLumpur.jpg} % Chapter heading image

\chapter{Legal}

%----------------------------------------------------------------------------------------
%	BIBLIOGRAPHY
%----------------------------------------------------------------------------------------

\chapterimage{chapter_head_9_Milan.jpg} % Chapter heading image

\addcontentsline{toc}{chapter}{\textcolor{trolleygrey}{References}}
\nocite{*}
\printbibliography[title={References}]
%------------------------------------------------


%----------------------------------------------------------------------------------------
%	COPYRIGHT PAGE
%----------------------------------------------------------------------------------------

\newpage
~\vfill
\thispagestyle{empty}

\noindent Copyright \copyright\ 2018 Ripa Founder Team\\ % Copyright notice

\noindent \textsc{Pubblicato da Ripa Founder Team}\\ % Publisher

\noindent \textsc{www.ripaex.io}\\ % URL

\noindent Materiale protetto dalla Licenza MIT (di seguito la ``Licenza''). Non dovrai usare questo file eccetto in conformità con la Licenza. Puoi ottenere una copia della Licenza a \url{https://opensource.org/licenses/MIT}. A meno che non sia richiesto dalla lagislazione applicata o concordato via scritto, il software distribuito con la Licenza è distribuito \textsc{``così com'è'' senza garanzie di alcun tipo}, esplicite o implicite. Fate riferimento alla Licenza per il linguaggio specifico che governa i permessi e le limitazione indicati dalla Licenza.\\ 

\noindent \textit{Prima pubblicazione, Marzo 2018} % Printing/edition date

\end{document}
